%%Packages
\usepackage{latexsym, amsthm, amscd, amsfonts, mathrsfs, amsmath, amssymb, stmaryrd, tikz-cd, mathrsfs, bbm, url, esint, mathtools, enumerate, caption, subcaption}
\usepackage{dsfont}
\usepackage{blkarray}
\usepackage[utf8]{inputenc} % allow utf-8 input
\usepackage[T1]{fontenc}    % use 8-bit T1 fonts
\usepackage{hyperref}       % hyperlinks
\usepackage{url}            % simple URL typesetting
\usepackage{booktabs}       % professional-quality tables
\usepackage{nicefrac}       % compact symbols for 1/2, etc.
\usepackage{microtype}      % microtypography
\usepackage{xcolor}  
%%algorithm
\usepackage{algorithm}
\usepackage{algpseudocode}
\usepackage{enumitem}
\usepackage{comment}






%%Commands and operators
\newcommand{\tr}{\mathrm{Tr}}
\newcommand{\He}{\mathrm{He}}
\newcommand{\Corr}{\mathrm{corr}}
\newcommand{\cmin}{c_{\min}}
\newcommand{\Cov}{\mathrm{cov}}
\newcommand{\diag}{\mathrm{diag}}
\newcommand{\pr}{\mathbb{P}}
\newcommand{\E}{\mathbb{E}}
\DeclareMathOperator{\CP}{CP}
\DeclareMathOperator{\Var}{Var}
\newcommand{\Prr}{\mathop{\rm Pr}\nolimits}
\DeclareMathOperator{\argmin}{argmin}
\DeclareMathOperator{\argmax}{argmax}
\DeclareMathOperator{\Maj}{Maj}
\DeclareMathOperator{\Ramp}{Ramp}
\DeclareMathOperator{\ReLU}{ReLU}
\DeclareMathOperator{\Unif}{Unif}
\DeclareMathOperator{\Rad}{Rad}
\DeclareMathOperator{\INAL}{INAL}
\DeclareMathOperator{\sign}{sgn}
\DeclareMathOperator{\erf}{erf}
\DeclareMathOperator{\erfc}{erfc}
\DeclareMathOperator{\orb}{orb}
\DeclareMathOperator{\Span}{span}
\DeclareMathOperator{\NS}{NS}
\DeclareMathOperator{\NN}{NN}
\DeclareMathOperator{\SBM}{SBM}
\DeclareMathOperator{\poly}{poly}
\DeclareMathOperator{\TV}{TV}
\DeclareMathOperator{\Inf}{Inf}
\newcommand{\cF}{\mathcal{F}}
\newcommand{\cX}{\mathcal{X}}
\newcommand{\cU}{\mathcal{U}}
\newcommand{\cD}{\mathcal{D}}
\newcommand{\cN}{\mathcal{N}}
\newcommand{\cE}{\mathcal{E}}
\newcommand{\cP}{\mathcal{P}}
\newcommand{\cG}{\mathcal{G}}
\newcommand{\cI}{\mathcal{I}}
\newcommand{\bR}{\mathbb{R}}
\newcommand{\bI}{\mathbb{I}}
\newcommand{\bN}{\mathbb{N}}
\newcommand{\cK}{\mathcal{K}}
\newcommand{\cC}{\mathcal{C}}
\newcommand{\cL}{\mathcal{L}}
\newcommand{\cB}{\mathcal{B}}
\newcommand{\bS}{\mathbb{S}}
\newcommand{\w}{\mathbf{w}}
\newcommand{\x}{\mathbf{x}}
\newcommand{\indep}{\perp\!\!\!\perp}




%%Colors
\definecolor{mydarkblue}{rgb}{0,0.08,0.45}
  \hypersetup{ %
    colorlinks=true,
    linkcolor=mydarkblue,
    citecolor=mydarkblue,
    filecolor=mydarkblue,
    urlcolor=mydarkblue,
    }
\definecolor{DSgray}{cmyk}{0,0,0,0.7}
\definecolor{DSred}{cmyk}{0,0.7,0,0.7}


%%Author notes
\newcommand{\Authornote}[2]{\noindent{\small\textcolor{DSgray}{\sf{
\textcolor{purple}{[#1: #2]}}}}}
\newcommand{\Authornoteo}[2]{\noindent{\small\textcolor{DSgray}{\sf{
\textcolor{orange}{[#1: #2]\marginpar{\textcolor{orange}{\fbox{\Large !}}}}}}}}
\newcommand{\Authornotec}[2]{\noindent{\small\textcolor{DSgray}{\sf{
\textcolor{blue}{[#1: #2]\marginpar{\textcolor{blue}{\fbox{\Large !}}}}}}}}
\newcommand{\enote}{\Authornote{Elisabetta}}
\newcommand{\anote}{\Authornote{Aryo}}


%%Theorem styles
\theoremstyle{plain}
\newtheorem{definition}{Definition}%[section]
\theoremstyle{plain}
\newtheorem{assumption}{Assumption}%[section]
\newtheorem{ass}{Assumption}
\theoremstyle{plain}
\newtheorem{theorem}{Theorem}%[section]
\theoremstyle{plain}
\newtheorem{informaltheorem}{Informal Theorem}%[section]
\theoremstyle{plain}
\newtheorem{exe}{Exercise}%[section]
\theoremstyle{plain}
\newtheorem{example}{Example}%[section]
\theoremstyle{plain}
\newtheorem{fact}{Fact}%[section]
\theoremstyle{plain}
\newtheorem{proposition}{Proposition}%[section]
\theoremstyle{plain}
\newtheorem{claim}{Claim}%[section]
\theoremstyle{plain}
\newtheorem{lemma}{Lemma}%[section]
\theoremstyle{plain}
\newtheorem{corollary}{Corollary}%[thm]
\theoremstyle{plain}
\newtheorem{conjecture}{Conjecture}%[section]
\theoremstyle{remark}
\newtheorem{remark}{Remark}%[section]
\theoremstyle{remark}
\newtheorem{note}{Note}%[section]
% \theoremstyle{discussion}
\newtheorem{discussion}{Discussion}%[section]
\theoremstyle{plain}
\newtheorem{conj}{Conjecture}%[section]